\section*{Scope, sources and verification status}
\addcontentsline{toc}{section}{Scope, sources and verification status}
This is a consolidated working exposition, not a completed unconditional
proof of RH. The prime-only analytic ending and signed finite comparison
are printed first. The elementary arithmetic packet, minus and dyadic
bounds, global high-twist bound and common-integer finiteness arguments
are proved next. Both uniform character estimates are proved in
Theorem~\ref{thm:sr-generic-character-estimate} and
Proposition~\ref{prop:sr-nontrivial-final-assembly}.
Corollary~\ref{cor:sr-principal-plus-height-proved} assembles
their polynomial plus-height bound.

The deformation argument now includes actual classical ordinary seeds,
their same-module unitary insertion, the cyclotomic arithmetic frame,
its uniform class-rank and strict-cohomology bounds, and the actual entry
space with flat charts, selected curves and sharp ideal.
The rank-four tower witnesses and exact descent are proved in
Section~\ref{sec:sr-rank-four-actual-endpoint}.
Section~\ref{sec:sr-rank-three-host-endpoints} gives finite host algebras
and exact automorphic endpoints from the stated common support output.
Section~\ref{sec:actual-seed-two-propagations} proves the common
rank-two/rank-three propagation output. The actual local crystalline
component and original-level comparison
are proved, including the nearest-weight case. The rank-four common space, good residual lattice, actual modular
ideal, residual paths and nonsplit block extension are constructed.
Its final finite-weight and no-block arithmetic endpoint is proved
in Proposition~\ref{prop:sr-four-actual-final-endpoint}.
Proposition~\ref{prop:sr-gaussian-actual-g10} specializes the
same rank-two/rank-three construction to the actual Gaussian
frame, including coefficient prime three and the boundary weight. No uniform character estimate is inferred
from a seed alone.

The geometric chapters construct the archimedean common Fock quotient
and all finite local matching maps,
including inert places in either rank parity, and prove the exact
good-place coinvariant condition.
Subsection~\ref{subsec:sr-theta-decay-closed} proves both ordinary
convergence and the small-side rapid-decay and moderate-growth clauses
with the specified bar-model test.
The canonical cusp-weight mean has constant four and the quadratic
CM period is evaluated. Section~\ref{sec:sr-compact-global-comparison}
also proves the nonzero global comparison at the compact base.
The finite cone vacuum, absolutely convergent adelic Euler
factorization and four global constant-term producers are also proved.
The all-signature global comparison, quotient descent, rank recursions
and the general integral-model, saturation and covolume bound
remain to be completed. The universal principal Lie-bundle
slope estimate is proved in
Theorem~\ref{thm:sr-general-principal-slope}.

The term \emph{connected} in the formal family predicate means one
compatible eventual window witness, as explained in
Section~\ref{sec:sr-latest-community-check}.
The required local inequalities and quantifier order are retained in
Definition~\ref{def:sr-supply}.
The second part preserves the earlier translations of the longer route.
Their explicit hypotheses remain in force, with the newer first-part
proofs discharging only the clauses identified in the closing ledger.

The shorter route follows Anton's Principia files
\href{https://principia-math.com/projects/d16ab288ed353bc5b10e292ace0ddbc1/editor?file=d59903a6498c7db021b728233ee6bdc2}{\srcid{FourthSamplesRH}},
\href{https://principia-math.com/projects/d16ab288ed353bc5b10e292ace0ddbc1/editor?file=dfeee1c759dca3c2b8e378a2dec2c780}{\srcid{PrimeFourthSamples}},
\href{https://principia-math.com/projects/d16ab288ed353bc5b10e292ace0ddbc1/editor?file=57674e7030ca8ce9bc2b4113ee37ab82}{\srcid{SignedSignal}},
\href{https://principia-math.com/projects/d16ab288ed353bc5b10e292ace0ddbc1/editor?file=150a1ccfbf0ccf9a0b4297260c71c05d}{\srcid{PointwiseGeometry}},
\href{https://principia-math.com/projects/d16ab288ed353bc5b10e292ace0ddbc1/editor?file=116c16d3654fa8b487602e67ead50b96}{\srcid{SignedGeometry}},
and
\href{https://principia-math.com/projects/d16ab288ed353bc5b10e292ace0ddbc1/editor?file=41bb998820fe961dbf8e023c9469578b}{\srcid{ArchiveThin}}.
The separate
\href{https://principia-math.com/projects/d16ab288ed353bc5b10e292ace0ddbc1/editor?file=c36fb16dbf85efd1520d208bf47edf3a}{\srcid{PrincipalWindows}}
interface isolates the conductor-one connected-family route.
Its explicit geometric premises must still be proved in a standalone text.
The 552-line specialization, its 630-line cubic-logarithmic revision,
the subsequent two-character height reduction, and the historical
five-PDF review are recorded in
Section~\ref{sec:sr-source-review}; the implications it proves have
not been mistaken for proofs of those premises.

The
\href{https://github.com/leanprover-community/riemann-hypothesis/pull/12}{signed-signal upstream supplement}
at commit \srcid{9430017510358b5ed1ab12c8bb900fd706e60a32}
has an independent successful Lean 4.33.1 source rebuild of its 41 project
dependencies. That verifies the supplement's implications with their
stated hypotheses. It does not verify this new prose edition or prove
its supply assertion.
The simplification of the real Mellin identity by
\href{https://principia-math.com/projects/d16ab288ed353bc5b10e292ace0ddbc1/editor?file=58b04d0292b512962984e43972698e31}{uimspeedrun}
preserves its statement, deriving it from the complex identity.
We use the same specialization principle, with the actual integral
calculation printed in the text.

Recent geometric translations by Baelson, reviewed by lowrank, include
\href{https://principia-math.com/api/rh/step?name=Development.AbelianVarietyTheory.Proofs.DivisorProofs.ArCalcOfPrincipalisation}{arithmetic intersection calculus from principalisation},
\href{https://principia-math.com/api/rh/step?name=Development.AbelianVarietyTheory.Proofs.DivisorProofs.SemiPullOfPrincipalisation}{semipositivity under pullback},
\href{https://principia-math.com/api/rh/step?name=Development.AbelianVarietyTheory.Proofs.DivisorProofs.RigidBoundaryChartsForFanUnitaryDatumStatement}{rigid toroidal boundary charts}, and
\href{https://principia-math.com/api/rh/step?name=Development.AbelianVarietyTheory.Proofs.DivisorProofs.ExistsFramesFockMapPolynomialLieActionsStatement}{Fock frames and Lie actions}.
They locate the required geometric dependencies.
The short re-export modules do not themselves contain the proofs of
all their imported results.

The ordinary-zeta oscillation argument in this edition excludes positive
real zeros with the alternating zeta series; it consequently omits the
general real-pole case in the Lean argument.
The elementary bounds here also give the constant four in
\eqref{eq:sr-Elip}. These prose variations, including the dyadic pro-\(2\) group proof and
spectral uniqueness, dyadic matrix and integral-homology arguments,
have not been
separately formalized or checked by Lean.
The class-number proof prints the one-ramified-prime argument and the
strict cyclotomic descent explicitly, using the classical background
listed in Section~\ref{sec:sr-arithmetic}.

The full formal export has also been typeset as an indexed record edition.
That is a record transcription, not a human exposition.
Its reverse-conversion audit establishes fidelity to the input records,
not completion of this paper.
The
\href{https://github.com/ChristianKudeba/alpoge-rh-verification/releases/tag/rh-literal-audit-v1}{record edition and audit}
and
\href{https://math.mit.edu/~drew/grh/}{independent formal verification report}
are separate evidence from the mathematical arguments printed here.


\subsection*{The logical route through the paper}
The analytic target is Theorem~\ref{thm:sr-prime}: a
fourth-power prime-sampling estimate implies the Riemann
hypothesis for ordinary zeta. Theorem~\ref{thm:sr-window}
permits that estimate to be supplied on the admissible
windows. The signed finite comparison in
Theorem~\ref{thm:sr-supply} reduces it to the arithmetic
and geometric witnesses in Definition~\ref{def:sr-supply}.
These are the two substantive branches of the argument.

On the arithmetic branch, the minus and dyadic parts are
already bounded. The plus part is reduced, with its prime
weights and quantifiers intact, to
Definition~\ref{def:sr-character-inputs}. A single exponent
must bound the sum over every finite generic prime set;
a second exponent must bound each nontrivial exceptional
or resonant character. Both exponents precede the varying
host, character order and twist.
Finiteness for one host is useful but does not establish
these uniform estimates.

The deformation chapters build the mechanism for those
estimates: actual ordinary seeds, an arithmetic frame,
curve and link presentations, transport in one ordinary
module, and automorphic descent. The class-group constant
in Section~\ref{sec:arithmetic-frame-uniform-depth} is
chosen before cyclotomic depth. It need not itself be
polynomial in the host. Its role is to make the auxiliary
frame sufficiently deep; the final numerical character
count must be made after descent, with its stated uniform
conductor and weight bounds.

On the geometric branch, one must choose compatible data
on the same windows, prove the global theta comparison
and its rank recursions, and control the integral lattices
and metrics of those data. The local maps, convergence
estimates and fixed-model identities in the theta chapters
are components of this construction. They do not separately
choose the compatible family or prove the uniform
covolume bound. Here \emph{connected} refers to the
source's compatible eventual window witness, as explained
in Section~\ref{sec:sr-latest-community-check}.

The closing ledger in Section~\ref{sec:sr-source-review}
records the exact frontier in these two branches.
Until both supply obligations there have been discharged,
the final RH conclusion remains conditional in this prose
edition. Part II preserves the earlier longer-route
translations with their hypotheses and makes their
supporting arguments available to the first part.
